\uuid{CLWM}
\titre{Quotient/Composée/DL usuels}
\auteur{hanani ; mzari}
\organisation{exo7}
\licence{CC BY-NC-SA 4.0}
\datecreate{2020-01-20}
\niveau{L1}
\module{Analyse}
\chapitre{Développement limité}
\sousChapitre{Calculs}
\difficulte{3}
\theme{Opérations sur les DL, Quotient, Composée, DL usuels}
\exo7qcmid{lille-2.developpements-limites.013}

\contenu{
\texte{
Soit $a\in \Rr$ et $f$ une fonction. On dit que $f(x)$ est un $o\left((x-a)^n\right)$ au voisinage de $a$ et on écrit $f(x)=o((x-a)^n)$ si $\displaystyle \lim _{x\to a}\frac{f(x)}{(x-a)^n}=0$. On note par $DL_n(a)f(x)$ le développement limité de $f$ à l'ordre $n$ en $a$. On note aussi $DL_n(a^+)f(x)$ (resp. $DL_n(a^-)f(x)$) le développement limité de $f$ à l'ordre $n$ à droite en $a$ (resp. à gauche en $a$).
}

\begin{question}
\qclassification{125.02}
\qtags{motcle=Quotient/Composée/DL usuels}
Parmi les affirmations suivantes, cocher celles qui sont vraies :
\begin{answers}
\bad{Pour obtenir le $DL_2(0)\sqrt{2+t}$, on écrit :
$$\displaystyle \sqrt{2+t}=\sqrt{1+(1+t)}=1+\frac{(1+t)}{2}-\frac{(1+t)^2}{8}+o((1+t)^2).$$}
\bad{Pour obtenir le $DL_2(2)\sqrt{x}$, on écrit :
$$\displaystyle \sqrt{x}=\sqrt{1+(x-1)}=1+\frac{(x-1)}{2}-\frac{(x-1)^2}{8}+o((x-1)^2).$$}
\good{$\displaystyle DL_2(0)\frac{\ln (1+t)}{1+t}=t-\frac{3t^2}{2}+o(t^2)$.}
\good{$\displaystyle DL_2(1)\frac{\ln x}{x}=(x-1)-\frac{3(x-1)^2}{2}+o((x-2)^2)$.}
\end{answers}
\vskip2mm
\begin{explanations}
Pour obtenir le $DL_2(0)\sqrt{2+t}$, on écrit :
$$\displaystyle \sqrt{2+t}=\sqrt{2}\sqrt{1+\frac{t}{2}}=\sqrt{2}\left(1+\frac{t}{4}-\frac{t^2}{32}+o(t^2)\right).$$
Pour obtenir le $DL_2(2)\sqrt{x}$, on pose $t=x-2$. D'où
$$\sqrt{x}=\sqrt{2+t}=\sqrt{2}\left[1+\frac{(x-2)}{4}-\frac{(x-2)^2}{32}\right]+o((x-2)^2).$$
On a : $\displaystyle \frac{\ln (1+t)}{1+t}=\left(t-\frac{t^2}{2}+o(t^2)\right)\left(1-t+t^2+o(t^2)\right)=t-\frac{3t^2}{2}+o(t^2)$.
Enfin, en posant $t=x-1$, on obtient :
$$\frac{\ln x}{x}=\frac{\ln (1+t)}{1+t}=t-\frac{3t^2}{2}+o(t^2)=(x-1)-\frac{3(x-1)^2}{2}+o((x-1)^2).$$
\end{explanations}
\end{question}
}
